\dykband{EXERCISES IN SWAFL}
{\small Open SWAFL (desktop or online) with the benchmark project. The chapter of the \href{https://github.com/halilunaziru73-creator/Coupled-Water-Fluxes-Pore-Scale-Dynamics-Precision-Irrigation-Mulching-Feedback-Loops/raw/main/Software/v1.1.0/SWAFL_Desktop/docs/manual_v1.1.0/SWAFL_Manual_v1.1.0.pdf}{SWAFL manual} that explains each task is given in brackets; record your values in a table and compare them with what is expected (\emph{Check}).}\par\vspace{3pt}
\begin{exlist}
\item {\color{khblue}[Manual, Chapter 11]} Open \emph{Surface irrigation design} in \emph{Your soil} mode for the loam. Record the intake, the advance time, the cut-off and the efficiencies.\par\noindent{\small\itshape Check: the design table lists them.}
\item {\color{khblue}[Manual, Chapter 11]} Use \emph{Best inflow}; compare the efficiency at the best inflow with half and double that inflow.\par\noindent{\small\itshape Check: the efficiency is highest near the best inflow.}
\item {\color{khblue}[Manual, Chapter 11]} Read the tornado chart of the sensitivity. Which input changes the efficiency most?\par\noindent{\small\itshape Check: usually the intake or the inflow.}
\item {\color{khblue}[Manual, Chapter 11]} Read the irrigation schedule and the message \emph{When to irrigate}. Check the next date against the sensor and the trigger.\par\noindent{\small\itshape Check: the date follows the depletion.}
\end{exlist}
